
\clearpage
\VerdierPhysicalPageAnchor{0125}%
\begingroup
\small
\setlength{\emergencystretch}{1.5em}%
\setlength{\abovedisplayskip}{4pt plus 2pt minus 2pt}%
\setlength{\belowdisplayskip}{4pt plus 2pt minus 2pt}%

\VerdierSemanticAnchor{ch2-1-3-5}
\noindent\textbf{命题 1.3.5.}\quad
\emph{设 $C$ 为一个阿贝尔范畴。函子}

\VerdierSemanticAnchor{ch2-1-3-5-h0-functor}
\[
  H^0\colon K(C)\longrightarrow C
\]
\noindent\emph{是一个上同调函子。}

\VerdierSemanticAnchor{ch2-1-3-5-first-exact-sequence}
事实上，由（第一章，3.3.1），只需证明对于 $\operatorname{tr}_1(C)$
（第一章，3.1.1）的每个对象 $X\xrightarrow{f}Y$，序列
\[
  H^0(X)\xrightarrow{H^0(f)}H^0(Y)
  \xrightarrow{H^0(p(f))}H^0(c(f))
  \xrightarrow{H^0(q(f))}H^0(X[1])
\]
\noindent 是正合的。

\VerdierSemanticAnchor{ch2-1-3-5-trii-reduction}
利用性质 TRII (1.1.1)，可见只需证明序列

\VerdierSemanticAnchor{ch2-1-3-5-reduced-exact-sequence}
\[
  H^0(Y)\xrightarrow{H^0(p(f))}H^0(c(f))
  \xrightarrow{H^0(q(f))}H^0(X[1])
\]
\noindent 是正合的；这立即由定义（第一章，3.1.2 与 3.2.2）得到。

\VerdierSemanticAnchor{ch2-1-3-5-proof-complete}
函子 $H^0\colon K(C)\longrightarrow C$ 称为\emph{典范上同调函子}。

\VerdierSemanticAnchor{ch2-1-3-5-shifted-cohomology}
按照惯例，令 $H^n(X)=H^0(X[n])$。

\VerdierSemanticAnchor{ch2-1-3-6}
\medskip
\noindent\textbf{命题 1.3.6.}\quad
\emph{设 $C$ 为一个阿贝尔范畴。下列条件等价：}

\VerdierSemanticAnchor{ch2-1-3-6-equivalent-conditions}
\begin{enumerate}
\renewcommand{\labelenumi}{\roman{enumi})}
\item \emph{范畴 $K(C)$ 是阿贝尔范畴。}
\item \emph{范畴 $K(C)$ 可分解 (1.2.8)。}
\item \emph{范畴 $C$ 可分解。}
\end{enumerate}

\VerdierSemanticAnchor{ch2-1-3-6-i-iff-ii}
等价关系 $\mathrm{i)}\Longleftrightarrow\mathrm{ii)}$ 由 (1.2.9) 得到。

\VerdierSemanticAnchor{ch2-1-3-6-ii-implies-iii}
由于 $C$ 是 $K(C)$ 的满子范畴，蕴含
$\mathrm{ii)}\Longrightarrow\mathrm{iii)}$ 是显然的。

\VerdierSemanticAnchor{ch2-1-3-6-hstar-complex}
对于 $C$ 中每个复形 $X$，用 $H^*(X)$ 表示复形
\[
  H^*(X)^l=H^l(X),\qquad l\in\mathbb Z,
\]
\[
  d_{H^*(X)}^l=0,\qquad l\in\mathbb Z.
\]

\VerdierSemanticAnchor{ch2-1-3-6-hstar-equivalence}
可以验证：当 $C$ 可分解时，函子
$H^*\colon K(C)\longrightarrow K(C)$ 是范畴等价。

\VerdierSemanticAnchor{ch2-1-3-6-proof-complete}
由此得到 $\mathrm{iii)}\Longrightarrow\mathrm{ii)}$。
\endgroup


